% Owner: corresponding author (collect input from everyone)
% Required. Does not count toward the page limit. At most 1 page.
% Policy: https://www.iclr.cc/Conferences/2027/AIPolicyForAuthors
% Confirm every item below with the team; the submission form must match this text.
\subsection*{AI use statement}

Generative AI is part of the method studied in this paper, and we also used it as a research and writing aid. We list both uses following the ICLR 2027 AI Policy for Authors.

\textbf{Uses with required disclosure.}
We used generative AI to generate synthetic datasets: Gemini-3.1-Pro-Preview proposes task instructions from explored application states and executes them to produce the candidate trajectories in OpenMobile-Data (Section~\ref{sec:data-pipeline}).
We used generative AI to clean datasets: on the Android emulator, a vision-language judge assesses whether each rollout completes its task, and rollouts it rejects are removed. Trajectories collected in \textsc{MobileGym++} are instead verified by programmatic checkers.
We used AI coding assistants to implement methods: the 35 new \textsc{MobileGym++} applications, their app-native tools, the task checkers, and parts of the training, evaluation, and data-collection pipelines. For each application, the authors first specified the data schema and interaction graph, and the assistants implemented the application against this specification (Section~\ref{sec:mgpp-apps}).
We did not use generative AI to develop theoretical models, formulate or prove mathematical claims, propose hypotheses, or design the research methodology and experiments. We also did not use it to translate text, perform qualitative analysis, or interpret results.

\textbf{Uses with recommended disclosure.}
We used AI assistants to edit and debug software code, to write plotting code for figures, to draft and edit parts of the paper for clarity and structure, and to format references.

\textbf{Verification.}
Human annotators verified every \textsc{MobileGym++} application against the corresponding real client on a device. App-native tools were tested by executing each call and checking the resulting change in application state, and task checkers were validated with test cases and manual spot checks. All synthesized trajectories pass completion checks before release, and the authors manually inspected samples of the released data. All experimental numbers come from our own runs, and all citations were checked against the original sources. We have reviewed all AI-assisted work and take responsibility for the final content of this paper, including text, claims, and artifacts produced with the aid of generative AI.
